\documentclass[11pt]{article}
\usepackage[margin=.82in]{geometry}
\usepackage{amsmath,amssymb,booktabs,hyperref}
\hypersetup{colorlinks=true,urlcolor=blue}
\newcommand{\F}{\mathbb F}
\title{\vspace{-1.6em}Exact root counts for two to five terms\\ in deep \(p\)-adic disks}
\author{A supplement to the ProveIt GRH article}
\date{27 November 2026}
\begin{document}
\maketitle
Let \(v_p(p)=1\). Fix a nonzero center \(A\) and the closed disk
\(v_p(X/A-1)\ge r\), with \(r>1/(p-1)\). Let \(N_p(t,r)\) be the greatest
number of distinct roots there of a nonzero polynomial with at most
\(t\) terms, allowing finite extensions of \(\mathbb Q_p\).
For \(2\le t\le \min(5,p)\) the exact answer is:
\begin{center}
\begin{tabular}{cl}
\toprule
\(t\) & \(N_p(t,r)\)\\
\midrule
2 & \(1\)\\
3 or 4 & \(p\) for \(r\le1/(p-t+1)\); \(t-1\) otherwise\\
5 & \(p+1\) for \(r\le1/(p-3)\);
 \(p\) for \(1/(p-3)<r\le1/(p-4)\); \(4\) otherwise\\
\bottomrule
\end{tabular}
\end{center}
The upper bounds hold even counting multiplicity. In every region
there are witnesses about \(A=1\) with integer coefficients,
which can also be chosen irreducible over \(\mathbb Q\).
For all \(2\le t\le p\) we further have
\(N_p(t,r)\le\max(t-1,p(t-2))\), uniformly in the depth.
For example, at \(p=t=5,\ r=1/4+10^{-9}\) the
factorial-only index bound of the article is at least \(488{,}281{,}250\);
the exact answer is six.

\section{Reduction of the moment recurrence}
The \href{https://github.com/VladimirReshetnikov/ProveIt/blob/dd990b27a22290c756c1bd4adce1d971bf7e7821/Analysis/GeneralizedRiemannHypothesis/Article/sections/smd.tex#L57}{article's local lemma}
sets \(X=A\exp(bY)\) with \(v_p(b)=r\), an isometric parametrization
after scaling. The coefficients in \(Y\) are \(M_k b^k/k!\), with
\(M_k=\sum_j c_j A^{n_j}n_j^k\) for the original coefficients and exponents.
It suffices to use rational depths; for other depths the asserted
upper bounds follow by choosing a slightly smaller rational depth in
the same region, and the lower witnesses below work unchanged.
By Vandermonde and the integral monic recurrence with polynomial
\(\prod_j(T-n_j)\), divide all moments by one of least valuation among
the first \(t\). After division they define an integral functional
on the power basis of the recurrence algebra, with some early value
of valuation zero. Write
\[
 w(k)=rk-v_p(k!)=(r-1/(p-1))k+S_p(k)/(p-1).
\]
We can assume exactly \(t\) distinct exponents by adjoining zero
coefficients if needed. For \(k\ge t\), digit sum \(S_p(k)\ge t-1\)
implies \(w(k)>\max_{i<t}w(i)=w(t-1)\).
Also, \(w(i)-w(k)<1\) for \(i<t\le k\), so a normalized
moment of valuation at least one excludes that index from attaining
the minimum coefficient valuation.
\newpage
For the remaining moments, work in
\[
 \mathcal R=\F_p[T]/\prod_j(T-\bar n_j).
\]
If \(T^k\) is in the span of earlier \(T^\ell\)'s with
\(S_p(\ell)\le S_p(k)\), lift the span identity modulo \(p\) to
the integral recurrence algebra. If the \(k\)-th normalized moment
has valuation less than one, at least one of those earlier moments
then has no larger valuation. Since \(w(\ell)<w(k)\), the earlier
coefficient has strictly smaller valuation. This works over ramified
fields too: reduction is modulo the ideal \((p)\), not merely the
maximal ideal.

Put \(U=T^p,\ V=T-U\). On the factor at a residue \(a\), we have
\(U=a\) and \(V=T-a\), since its multiplicity is at most \(p\).
Hence \(T^{p^j}=U\) for \(j\ge1\).
The base-\(p\) digits \(a_j\) of a surviving index \(k\) with
\(S_p(k)\le t-2\) compress it to \(a_0+p\sum_{j\ge1}a_j\) with
the same digit sum and the same class in \(\mathcal R\). Any strict
compression excludes \(k\). The largest possible remaining index
is at most \(p(t-2)\). Strassmann's theorem gives the claimed
uniform bound, since \(w(k)\) tends to infinity.

\section{Small residue spans and the endpoint jumps}
For \(t=3\), the only compressed large candidate is \(p\). For \(t=4\)
the additional classes \(TU\) and \(U^2\) both lie in the span of
\(1,T,T^2,U\). Indeed, with multiplicities at most two,
\(2TU=T^2+U^2\). If there are at most two distinct residues, \(U^2\)
is affine in \(U\). With at least three residues there is at most
one repeat \(a\), in which case \((U-a)V=0\) gives
\(U^2=T^2+2aU-2aT\); when there is no repeat one takes \(V=0\).
If a multiplicity is at least three and there are two residues
\(a,b\), repeated \(a\) and singleton \(b\), then
\(TU=aT+b(U-a)\). With a single residue all powers of \(U\) are constants.

For \(t=5\), \(U^2\) lies in the span of \(1,T,T^2,U,TU\):
use \(V^2=0\), or at most two residues, or the sole remaining
profile \((3,1,1)\), where \((U-a)V=0\).
We must also eliminate \(T^2U,TU^2,U^3\) using
\(1,T,T^2,T^3,U,TU\). If \(V^2=0\) and there are at most three
residues, reduce \(U^3\) to degree two in \(U\).
Then \(T^3-U^3=3U^2V\), which eliminates both mixed classes.
If there is at most one repeat \(a\) (or \(V=0\)), use
\[
 T^3-U^3=3a^2V\quad\text{if }V^2=0,\qquad
 T^2U=aT^2+U^3-aU^2,\qquad TU^2=a^2T+U^3-a^2U.
\]
Here \(U^3\) was already eliminated for at most three residues;
for at least four residues \(V^2=0\) gives the first identity.
The only case left is profile \((3,2)\). Both \(U^2,U^3\) are
affine in \(U\). To express \(T^2U\), choose a cubic in \(T\)
matching its first two derivatives at the triple residue and first
derivative at the double residue. This is quadratic Hermite
interpolation on its derivative, and integration is possible as
2 and 3 are units. Correct its two values by an affine function of \(U\).
This proves the spans. Every span uses earlier exponents of at most
the same digit sum, so the largest index left for five terms is \(p+1\).
\newpage
When \(r>1/(p-3)\) the index \(p+1\) drops out since
\(w(p+1)-w(4)=(p-3)r-1>0\).
For the final drop to \(t-1\), if \(k\ge p\) then
\[
 v_p(k!)\le \frac{k-1}{p-1}\le\frac{k-t+1}{p-t+1}.
\]
When \(t\le k<p\) the factorial valuation is zero. Thus
\(w(k)>w(t-1)\) for every \(k\ge t\) as soon as \(r>1/(p-t+1)\).
These estimates prove all the upper bounds.

\section{Witnesses and irreducibility}
For \(3\le t\le 5,\ p\ge t\), the following polynomial has \(p\)
distinct roots in the disk at depth \(1/(p-t+1)\):
\[
 X^p-\sum_{j=0}^{t-2}\binom pj(X-1)^j+p^{2t-1}.
\]
It has at most \(t\) terms. In \(Z=X-1\) the lower Newton polygon
has slopes \(-2\) and \(-1/d\), where \(d=p-t+1\):
the vertices are \((0,2t-1),(t-1,1),(p,0)\).
At depth 2 the initial binomial has \(t-1\) simple nonzero roots
over the algebraic residue field, and at depth \(1/d\) the
normalized reduction has form \(Y^{t-1}(c+Y^d)\) with \(c\ne0\)
and \(d\) simple nonzero roots. Lift both by Hensel.

To attain \(p+1\) for five terms, take
\[
 F(X)=X^{2p}+1-p^2(X^{p+1}+X^{p-1})+
                [2(p^2-1)+p^5]X^p.
\]
Divide out \(X^p\) and substitute \(X=\exp(bY)\) with \(v_p(b)=r\).
The constant is \(p^5\), odd coefficients and the quadratic vanish,
and every even \(k\ge4\) has valuation \(2+w(k)\):
its coefficient is \(2(p^k-p^2)b^k/k!\).
At \(r=1/(p-3)\), the only minimal indices are \(4,p+1\).
They tie; for even \(4<k<p\) the weights strictly increase,
and for even \(k>p+1\) the digit sum is even and at least two,
so \(w(k)>w(p+1)\). The normalized reduction is
\(Y^4(c+d'Y^{p-3})\), with \(c,d'\ne0\), giving \(p-3\)
simple nonzero roots to lift. At \(r=3/4\) the only minimal indices
are \(0,4\): \(2+w(4)=5\), and for even \(k\ge6\),
\(w(k)\ge(3/4-1/(p-1))k+2/(p-1)>3\).
This gives four more simple roots at the deeper depth.
The two root sets are disjoint by their valuations.

For the remaining \(t-1\) lower bounds take a monic polynomial
with \(t-1\) distinct integer roots arbitrarily close to 1.
All witnesses can be made Eisenstein at a prime different from \(p\)
without adding terms: keep the leading monomial, preserve every
other coefficient modulo a sufficiently high power of \(p\), and
impose the Eisenstein congruences by the Chinese remainder theorem.
The normalized reductions and their simple roots are preserved
under these perturbations. For the product of linear factors one
uses disjoint smaller disks about its simple roots and the same
argument there. Finally scale the variable for a general center
(the assertion of rational irreducibility is about center 1).

\paragraph{Scope and replay.}
The related \href{https://arxiv.org/abs/1008.4808}{Avenda\~no--Krick work}
bounds roots in a fixed local field. Here the disk-depth extrema
range over finite extensions. The accompanying exact script tests
the residue spans and Newton profiles independently. The proof
uses neither RH nor GRH and claims neither publication priority
nor a Lean formalization. It does not settle the \(-87\) case.
\end{document}
